\documentclass[10pt,a4paper]{amsart}
\usepackage[margin=19mm]{geometry}
\usepackage{amsmath,amssymb,mathtools}
\usepackage[scr=rsfs,cal=euler]{mathalfa}
\usepackage{xfrac}
\usepackage[unicode,hidelinks,bookmarks=false]{hyperref}
\newcommand{\ie}{i.e.\ }
\newcommand{\cf}{cf.\ }
\newcommand{\resp}{respectively\ }
\newcommand{\Subsection}{\subsection}
\providecommand{\nobreakdash}{\nobreak\hbox{-}}
\setlength{\emergencystretch}{2em}
\multlinegap0pt
\usepackage{bm}
\usepackage{bbm}
\usepackage{dsfont}
\usepackage{xspace}
\usepackage{cancel}
\usepackage{mathtools}
\usepackage{amsthm}
\makeatletter
\@ifundefined{prop}{\newtheorem{prop}{Proposition}}{}
\makeatother
\DeclarePairedDelimiter{\abs}{\lvert}{\rvert}
\newcommand{\R}{{\mathbb{R}}}
\newcommand{\cB}{{\mathcal{B}}}
\newcommand{\cE}{{\mathcal{E}}}
\newcommand{\cI}{{\mathcal{I}}}
\newcommand{\cR}{{\mathcal{R}}}
\newcommand{\cX}{{\mathcal{X}}}
\newcommand{\indi}[1]{\mathds{1}_{#1}}
\title[Factorization by extremal privacy mechanisms: new insights i]{Factorization by extremal privacy mechanisms: new insights into efficiency}
\author{Chiara Amorino, Arnaud Gloter}
\date{Source: arXiv:2507.21769v1 (CC BY 4.0)}
\begin{document}
\maketitle

\noindent\emph{Source and licence.} Factorization by extremal privacy mechanisms: new insights into efficiency. Version arXiv:2507.21769v1 (CC BY 4.0), distributed under the Creative Commons Attribution 4.0 International licence, as recorded in the arXiv licence metadata for this exact version and shown on the abstract page https://arxiv.org/abs/2507.21769v1. Full rights notice: licenses/arxiv-2507.21769-CC-BY-4.0.txt.\par\smallskip

\noindent\textit{Editorial scope.} Counted unit: the Prop whose source label is \texttt{prop : Fisher unif} in the article (source0.tex lines 1242--1248, source label \texttt{prop : Fisher unif}), delivered together with the article's own complete original proof of it (source0.tex lines 1249--1308, one \texttt{proof} environment of 60 lines). No mathematics is written, repaired or re-derived: statement and proof are the article's own text line for line. Required context retained uncounted: eq : def evaluation operator (lines 1645--1649); p: evaluation (lines 1687--1708). Every definition, display and label of those lines is retained unchanged. Omitted from this package and not redistributed: 1--1241, 1309--1644, 1650--1686, 1709--3289. Nothing inside a retained excerpt is omitted or altered: every retained body is delivered exactly as the article's lines are written, so no omission receipt of any kind accompanies this package. Citations: the retained excerpts carry no citation command at all, so no bibliography entry of the article is transcribed and no reference is redistributed. Reference adaptation: none was needed and none was made; every reference inside the delivered unit resolves inside it, so no unresolved reference or citation is produced. Printed slips kept literal: no printed word, symbol, hypothesis or number of the counted statement or of its proof was altered. Layout and class: the article's own class is \texttt{article} with packages including geometry, babel, inputenc, fontenc, lmodern, yhmath, comment, bm, enumerate, amsthm, amsmath, amssymb, amsbsy, bbm; the wrapper is the lane's pinned \texttt{amsart} editorial wrapper at 10pt on A4, which restates the theorem-like environments the retained text needs and adds the article's own macro definitions that the retained text needs (\texttt{\textbackslash R}, \texttt{\textbackslash abs}, \texttt{\textbackslash cB}, \texttt{\textbackslash cE}, \texttt{\textbackslash cI}, \texttt{\textbackslash cR}, \texttt{\textbackslash cX}, \texttt{\textbackslash indi}), with extra wrapper packages (bm, bbm, dsfont, xspace, cancel, mathtools). The delivered numbering is the unit's own. Assets: no retained span contains a figure environment, a graphics-inclusion command, a PDF-page inclusion, a table or a TikZ picture; no image file is copied into this package, the package declares no asset dependency and no image file is redistributed with it. Licence: version arXiv:2507.21769v1 of this article is distributed under the Creative Commons Attribution 4.0 International licence, as recorded in the arXiv licence metadata for this exact version and shown on the abstract page https://arxiv.org/abs/2507.21769v1. A full rights notice is delivered at licenses/arxiv-2507.21769-CC-BY-4.0.txt. Characters: every retained excerpt is pure ASCII, so the delivered engine, XeLaTeX with its default Latin Modern text font, reports no missing character.
\begin{align}\nonumber %\label{eq: evaluation}
e_x : \quad & \widehat{B} \rightarrow \mathbb{R} \\
\label{eq : def evaluation operator}
& b \mapsto \limsup_{\epsilon \to 0} \frac{1}{|B(x, \epsilon)|} \int_{B(x, \epsilon)\cap\cX} b(y) \, dy,
\end{align}

\begin{prop}{\label{p: evaluation}}
	1) The map 
	\begin{align*}
		&\widehat{B}\times \cX \to \R
		\\
		&(f,x) \mapsto e_x(f)
	\end{align*}
	is measurable with respect to the $\sigma$-field $\cB(\widehat{B}) \otimes \cB(\cX)$.
	
	2) Let $\mu$ be a sub-probability on the space $\widehat{B}$ and assume that $b=\int_{\widehat{B}} f \mu(df)$. Then,
	%	Let $B$ be the set in Equation \ref{eq: def B}. Let us consider $\mathcal{E} \subset B$ the set of extreme points of $B$. Then, for any $b \in B$ and 
	for almost every $x \in \cX$, 
	\begin{equation}\label{eq : evalutation rigorous}
		b(x) =e_x(b) =\int_{\widehat{B}} e_x(f) d\mu (f). 
	\end{equation}
	3) The $\mu \otimes dx$ measure of the complement of the following set is zero:
	\begin{equation} \label{eq : neg by Fubini}
		\big\{(f,x)\in\widehat{B}\times\cX : e_x(f)  = \lim_{\epsilon \to 0} \frac{1}{|B(x, \epsilon)|} \int_{B(x, \epsilon)\cap\cX} f(y) \, dy\big\}.
	\end{equation}
	
	
\end{prop}

\begin{prop}\label{prop : Fisher unif}
	 For almost every $\theta_0 \in (0,\infty)$ the model
	$\tilde{p}_\theta(dr)$ is DQM at $\theta_0$. The following upper bound on the Fisher information holds
	\begin{equation*}
		\cI_{\theta_0}\left((\widetilde{p}_\theta)_\theta\right) \le \frac{(e^\alpha-1)^2}{\theta_0^2}.
	\end{equation*}
\end{prop}

\begin{proof}
	The density of the model is given by 
	\begin{equation}
		\label{eq : tilde p unif}
		\tilde{p}_\theta(r) =\frac{\tilde{p}_\theta(dr)}{\mu(dr)}=\frac{1}{\theta} \int_0^\theta r(x)dx.
	\end{equation}
		 We have to prove that this density admits a derivative with respect to the parameter $\theta$. Define 
	\begin{equation*}
		\cR=\left\{(r,x)\in B \times [0,\infty) \mid e_x^+(r):=\lim_{h\to0} h^{-1}\int_x^{x+h} r(y) dy \text{ exists} \right\}.
	\end{equation*}
	By similar argument as for the proof of the measurability of the operator $e_x$ defined by \eqref{eq : def evaluation operator}, we can check that set $\cR$ is measurable with respect to $\cB(B)\otimes\cB([0,\infty))$ (see Step 4 in the proof of Proposition \ref{p: evaluation}). By Fubini theorem, we deduce
	$\int_{\cE \times [0,\infty)} \indi{\cR^c}(r,\theta) \mu(dr)\times d\theta=
\int_{\cE} \left(\int_0^\infty \indi{\cR^c}(r,\theta) d\theta \right)\mu(dr)=0$, as for any $r\in \cE$, $\lim_{h\to0} h^{-1}\int_\theta^{\theta+h} r(y) dy$ exists for almost every $\theta$ and is equal to $r(\theta)$. Now, a further application of Fubini theorem entails  
$ \int_0^\infty\left(\int_{\cE} \indi{\cR^c}(r,\theta) \mu(dr) \right)d\theta=0$.
 Consequently, there exists $\overline{\Theta}$ a set of real, with full Lebesgue
  measure, such that $\int_{\cE} \indi{\cR^c}(r,\theta) \mu(dr)=0$ for all $\theta \in \overline{\Theta}$.

Now, we consider any $\theta_0 \in \overline{\Theta}$. From the definition of $\cR$, we have that $\mu(dr)$ a.e., the function $\theta \mapsto \int_0^\theta r(x)dx$ admits a derivative at $\theta=\theta_0$, denoted as $e^+_{\theta_0}(r)$. Recalling \eqref{eq : tilde p unif}, we deduce that $\tilde{p}_\theta(r)$ admits a derivative at $\theta=\theta_0$ given by
\begin{align} \nonumber
	\frac{\partial}{\partial \theta} \tilde{p}_{\theta_0}(r)
	&= -\frac{1}{\theta_0^2}\int_0^{\theta_0} r(x)dx+\frac{e^+_{\theta_0}(r)}{\theta_0} 
	\\ \label{eq : p tilde dot unif}
	&= \frac{1}{\theta_0^2}\int_0^{\theta_0} (e^+_{\theta_0}(r)-r(x))dx.
\end{align}
We recall the definition of $F^+_r = r^{-1} \left( \{ e^\alpha \} \right),$
and that \( r(x) = 1 + (e^{\alpha} - 1) \, \indi{F^+_r}(x) \).  
Substituting this expression into the definition of \( e^+_{\theta_0}(r) \), we obtain  
\[
e^+_{\theta_0}(r) = 1 + (e^{\alpha} - 1) \, e^+_{\theta_0} \left( \indi{F^+_r} \right).
\]
It implies
$e^+_{\theta_0}(r)-r(x)=(e^{\alpha}-1)(e^+_{\theta_0}(\indi{F^+_r})-\indi{F^+_r}(x))$.
%From \eqref{eq : tilde p unif}--
Inserting in \eqref{eq : p tilde dot unif}, we deduce that
\begin{equation*}
	\frac{(\frac{\partial}{\partial \theta} \tilde{p}_{\theta_0}(r))^2}{\tilde{p}_{\theta_0}(r)}
	=  \frac{(e^\alpha-1)^2}{\theta_0^4}\frac{\left( \int_0^{\theta_0} 
	(e^+_{\theta_0}(\indi{F^+_r})-\indi{F^+_r}(x))dx    \right)^2}{\widetilde{p}_{\theta_0}(r)}.
\end{equation*}
Collecting with  \eqref{eq : tilde p unif}, we deduce the expression for the Fisher information
\begin{equation*}
	\cI_{\theta_0}\left((\widetilde{p}_\theta)_\theta\right) 
	=	\frac{(e^\alpha-1)^2}{\theta_0^3}
	\int_{\cE}
\frac{\left( \int_0^{\theta_0} 
		(e^+_{\theta_0}(\indi{F^+_r})-\indi{F^+_r}(x))dx    \right)^2}{ \int_0^{\theta_0} r(x)dx } \mu(dr).
\end{equation*}	
	Since $e^+_{\theta_0}(\indi{F^+_r}) \in [0,1]$ and $\indi{F^+_r}(x)\in \{0,1\}$, we deduce that 
	\begin{equation} \label{eq : F_max unif}
	\abs*{ \int_0^{\theta_0} 
	(e^+_{\theta_0}(\indi{F^+_r})-\indi{F^+_r}(x))dx} \le \theta_0. 
	\end{equation}
	From $r\ge1$, we have $\int_0^{\theta_0} r(x)dx \ge \theta_0$, which gives
\begin{equation*}
	\cI_{\theta_0}\left((\widetilde{p}_\theta)_\theta\right) 
	\le	\frac{(e^\alpha-1)^2}{\theta_0^3}
	\int_{\cE} \frac{\theta_0^2}{\theta_0} \mu(dr) \le \frac{(e^\alpha-1)^2}{\theta_0^2},
\end{equation*}	
	where we used that $\mu$ is a sub-probability.	
\end{proof}

\end{document}
